\documentclass[
    preprint,
%    twocolumn,
    linenumbers,
    amsmath,amssymb,amsfonts,
    aps, physrev,
    superscriptaddress,
%    floatfix,
%    draft,
    10pt,
    bibnotes
    ]{revtex4-2}

\usepackage{bm}
\usepackage[T1]{fontenc}
\usepackage{inputenc}
\setcounter{secnumdepth}{3}
\usepackage{upgreek,babel,calc,mathrsfs,textgreek,bbm}
\usepackage{amsmath,amssymb,mathtools,graphicx,xcolor}


\newcommand{\erf}{\text{erf}}
\newcommand{\mt}[1]{{\color{green}{[MT: #1]}}}
\newcommand{\ld}[1]{{\color{red}{[LD: #1]}}}

\bibliographystyle{apsrev4-2}

\usepackage[unicode=true,pdfusetitle,
bookmarks=true,bookmarksnumbered=false,bookmarksopen=false,
breaklinks=false,pdfborder={0 0 1},backref=false,colorlinks=true]
{hyperref}


\begin{document}

\title{Random generalized Lotka-Volterra model on sparse graphs:\\
Gaussian expansion of dynamic cavity method}

\author{Mattia Tarabolo}
\email{mattia.tarabolo@polito.it}
\affiliation{Politecnico di Torino, Torino, Italy}

\author{A. R. Batista-Tomás}
\affiliation{University of Havana, La Habana, Cuba}

\author{Roberto Mulet}
\affiliation{University of Havana, La Habana, Cuba}

\author{Luca Dall'Asta}
\affiliation{Politecnico di Torino, Torino, Italy}
\affiliation{Collegio Carlo Alberto, Torino, Italy}
\affiliation{Italian Institute for Genomic Medicine (IIGM), Candiolo, Italy}

\date{\today}

\begin{abstract}
An article usually includes an abstract, a concise summary of the
work covered at length in the main body of the article. 
\end{abstract}

\maketitle

\section{\label{sec:intro}Introduction}

Intro on the subject.

\section{Method}

\subsection{Random generalized Lotka-Volterra model on sparse graphs}

We consider a generalized Lotka-Volterra model describing the dynamics
of an interacting community of $N$ species, labeled by $i=1,\dots,N$.
The species interacts on an undirected random graph $G$ with symmetric
adjacency matrix $A=\{a_{ij}\}$, meaning that species $i$ and $j$
are interacting if $a_{ij}=a_{ji}=1$, otherwise the matrix elements
are zero. The time-dependent number density of individuals of species
$i$ is denoted by $x_{i}(t)$, and evolves in time according to

\begin{equation}
\frac{d}{dt}x_{i}(t)=x_{i}(t)\left[1-x_{i}(t)+\sum_{j\neq i}a_{ij}J_{ij}x_{j}(t)+\eta_{i}(t)\right],\label{eq:Lotka-Volterra_eq}
\end{equation}

where $\eta_{i}(t)$ is a Gaussian noise with mean and variance

\begin{equation}
\langle\eta_{i}(t)\rangle=0,\qquad\langle\eta_{i}(t)\eta_{i'}(t')\rangle=2T\delta_{ii'}\delta(t-t').
\end{equation}

The couplings $J_{ij}$ represent the (per capita) effect of species
on one another. A negative coefficient $J_{ij}$ indicates a competitive
effect of species $j$ on species $i$.

We consider here Gaussian $J_{ij}$ with average $\overline{J_{ij}}=\overline{J_{ji}}=m/K$,
variance $\text{Var}(J_{ij})=\text{Var}(J_{ji})=\sigma^{2}/K$ and
covariance $\text{Corr}(J_{ij},J_{ji})=\gamma$. The overbar $\overline{\dots}$
represents the average over the random couplings ensemble. The scaling
of the moments with the average degree of the graph $K$ is needed
to produce a well-defined fully connected limit $K\to\infty$. The
parameter $-1\leq\gamma\leq1$, being related to the correlation between
the interaction coefficients $J_{ij}$ and $J_{ji}$, measures the
fraction of predator-prey pairs. Such a predator-prey pair consists
of two species $i$ and $j$ for which $J_{ij}$ and $J_{ji}$ have
opposite signs. If for example $J_{ij}$ is positive and $J_{ji}$
is negative, the presence of species $i$ has a detrimental effect
on species $j$, while the presence of species $j$ is beneficial
for species $i$. For any pair $i<j$ we can write

\begin{equation}
J_{ij}=\frac{m}{K}+\frac{\sigma}{\sqrt{K}}z_{ij},\qquad J_{ji}=\frac{m}{K}+\frac{\sigma}{\sqrt{K}}z_{ji},
\end{equation}

where $z_{ij}$ and $z_{ji}$ are Gaussian random variables with $\overline{z_{ij}}=0$,
$\overline{z_{ij}^{2}}=1$, $\overline{z_{ij}z_{ji}}=\gamma$.

For $\gamma=1$ we have $z_{ij}=z_{ji}$ with probability one, for
$\gamma=0$ $z_{ij}$ and $z_{ji}$ are uncorrelated, and for $\gamma=-1$
we have $z_{ij}=z_{ji}$ with probability one. The latter, $\gamma=-1$,
corresponds to having all the interactions of the predator-prey type.
Increasing the value of $\gamma$ this fraction decrease, until the
opposite case , $\gamma=1$ of no predator-prey type interactions.
The other interactions can be of a mutualistic type ($J_{ij}$ and
$J_{ji}$ both positive) or strictly competitive ($J_{ij}$ and $J_{ji}$
both negative).

\subsection{Gaussian approximation of the dynamic cavity equations on a Random
Regular Graph (RRG)}

We focus here on the case of a Random Regular Graph (RRG) $G$, where
edges are randomly extracted with the constraint that each node has
exactly $K$ neighbors. 

We use the Gaussian Expansion Cavity method \cite{taraboloGaussianApproximationDynamic2024a}
to derive a Langevin equation for a representative species. This approach
begins by constructing the Martin-Siggia-Rose-Janssen-De Dominicis
(MSRJD) dynamic partition function through path integrals of the trajectories
\cite{martinStatisticalDynamicsClassical1973c,janssenLagrangeanClassicalField1976,dedominicisDynamicsSubstituteReplicas1978b}

\begin{equation}
Z=\left\langle \int\mathcal{D}\vec{x}\prod_{i}p_{i}\left(x_{i}\left(0\right)\right)\delta\left(\dot{x_{i}}(t)-x_{i}(t)\left[1-x_{i}(t)+\sum_{j\neq i}a_{ij}J_{ij}x_{j}(t)+\eta_{i}(t)+h_{i}(t)\right]\right)\right\rangle _{\vec{\eta}},
\end{equation}

where $\int\mathcal{D}\vec{x}=\prod_{i}\int\mathcal{D}x_{i}$ is a
functional integral over the trajectories of the system, and $h_{i}(t)$
is an external field, needed to compute response functions. Following
the steps of \cite{taraboloGaussianApproximationDynamic2024a} we
put forward a dynamic cavity Ansatz and expand the cavity marginals
for small couplings. This results in a set of effective de-coupled
Langevin equations on the cavity graph. For any edge $(i,j)$ on the
graph $G$ the effective cavity population size $x_{i\setminus j}(t)$
evolves under the following stochastic process
\begin{align}
\frac{d}{dt}x_{i\setminus j}(t) & =x_{i\setminus j}(t)\Biggl[1-x_{i\setminus j}(t)+\sum_{k\in\partial i\setminus j}J_{ik}\mu_{k\setminus i}(t)\nonumber \\
 & \quad+\sum_{k\in\partial i\setminus j}\int_{0}^{t}dt'J_{ik}R_{k\setminus i}(t,t')J_{ki}x_{i\setminus j}(t')+\xi_{i\setminus j}(t)+h_{i\setminus j}(t)\Biggr],\label{eq:eff_proc_cav}
\end{align}

where $h_{i\setminus j}(t)$ is an external field, needed to compute
response functions, and $\xi_{i\setminus j}(t)$ is a Gaussian colored
noise

\begin{equation}
\left\langle \xi_{i\setminus j}(t)\right\rangle _{i\setminus j}=0,\qquad\left\langle \xi_{i\setminus j}(t)\xi_{i\setminus j}(t')\right\rangle _{i\setminus j}=2T\delta(t-t')+\sum_{k\in\partial i\setminus j}J_{ik}^{2}C_{k\setminus i}(t,t').
\end{equation}

The terms $\mu_{k\setminus i}(t)$, $R_{k\setminus i}(t,t')$, and
$C_{k\setminus i}(t,t')$ are respectively the cavity mean abundance,
the cavity response function, and the cavity (connected) two-point
correlation function. They are obtained self-consistently from the
effective dynamics of the neighbors

\begin{equation}
\mu_{k\setminus i}(t)=\left\langle x_{k\setminus i}(t)\right\rangle _{k\setminus i},\quad R_{k\setminus i}(t,t')=\frac{\delta}{\delta h_{k}(t')}\left\langle x_{k\setminus i}(t)\right\rangle _{k\setminus i},\quad C(t,t')=\left\langle x_{k\setminus i}(t)x_{k\setminus i}(t')\right\rangle _{k\setminus i},
\end{equation}

where $\left\langle \dots\right\rangle _{k\setminus i}$ describes the average over the distribution of the effective process Eq. \eqref{eq:eff_proc_cav} for the edge $(k,i)$. The effective process is therefore de-coupled, and the interactions are captured by the self-consistent external fields $\mu_{k\setminus i}(t)$, $R_{k\setminus i}(t,t')$, and $C_{k\setminus i}(t,t')$. A detailed derivation can be found in Appendix \ref{sec:path_int_RRG}.

Once the cavity averages have been determined through self-consistency, one can obtain the effective de-coupled Langevin equations for the species population

\begin{align}
\frac{d}{dt}x_{i}(t) & =x_{i}(t)\Biggl[1-x_{i}(t)+\sum_{k\in\partial i}J_{ik}\mu_{k\setminus i}(t)\nonumber \\
 & \quad+\sum_{k\in\partial i}\int_{0}^{t}dt'J_{ik}R_{k\setminus i}(t,t')J_{ki}x_{i}(t')+\xi_{i}(t)+h_{i}(t)\Biggr],\label{eq:eff_proc_marg}
\end{align}

with Gaussian colored noise

\begin{equation}
\left\langle \xi_{i}(t)\right\rangle _{i}=0,\qquad\left\langle \xi_{i}(t)\xi_{i}(t')\right\rangle _{i}=2T\delta(t-t')+\sum_{k\in\partial i}J_{ik}^{2}C_{k\setminus i}(t,t').
\end{equation}

The average $\left\langle \dots\right\rangle _{i}$ is performed over the effective de-coupled process Eq. \eqref{eq:eff_proc_marg}.

\subsection{Naive method: Gaussian expansion}

The effective processes can be averaged over the disorder of both the RRG and the random couplings. Due to homogeneity of the system, this yields an effective process for a representative cavity species
abundance $x_{c}(t)$

\begin{align}
\frac{d}{dt}x_{c}(t) & =x_{c}(t)\Biggl[1-x_{c}(t)+\frac{K-1}{K}m\mu_{c}(t)\nonumber \\
 & \quad+\frac{K-1}{K}\left(\frac{m^{2}}{K}+\gamma\sigma^{2}\right)\int_{0}^{t}dt'R_{c}(t,t')x_{c}(t')+\xi_{c}(t)+h_{c}(t)\Biggr],\label{eq:eff_proc_cav_RRG}
\end{align}

where the Gaussian colored noise has mean and covariance

\[
\left\langle \xi_{c}(t)\right\rangle _{c}=0,\qquad\left\langle \xi_{c}(t)\xi_{c}(t')\right\rangle _{c}=2T\delta(t-t')+\frac{K-1}{K}\left(\frac{m^{2}}{K}+\sigma^{2}\right)C_{c}(t,t').
\]

The representative species abundance in the original graph evolves
under the following effective process

\begin{equation}
\frac{d}{dt}x(t)=x(t)\Biggl[1-x(t)+m\mu_{c}(t)+\left(\frac{m^{2}}{K}+\gamma\sigma^{2}\right)\int_{0}^{t}dt'R_{c}(t,t')x(t')+\xi(t)+h(t)\Biggr],\label{eq:eff_proc_marg_RRG}
\end{equation}

where the Gaussian colored noise has mean and covariance

\begin{equation}
\left\langle \xi(t)\right\rangle =0,\qquad\left\langle \xi(t)\xi(t')\right\rangle =2T\delta(t-t')+\left(\frac{m^{2}}{K}+\sigma^{2}\right)C_{c}(t,t').
\end{equation}

It seems that the averaging over disorder was done in a wrong way.
We show however that this corresponds to a Gaussian approximation
of the disorder averaged dynamical partition function. The calculations
can be found in Appendix \ref{app:naive_method}. The problem is that this approximation corresponds to a large connectivity expansion, already done by \cite{aguirre-lopezHeterogeneousMeanfieldAnalysis2024a}.

\subsubsection{\label{subsec:FP_analysis_RRG}Fixed point analysis}

In the limit $T\to0$ we can assume that the system reaches a steady
state which does not depends on the initial condition. We further
assume that the dynamics reaches a unique fixed-point $\lim_{t\to\infty}x_{c}(t)=x_{c}^{*}$,
$\lim_{t\to\infty}x(t)=x^{*}$, and therefore $\lim_{t\to\infty}\xi_{c}(t)=\xi_{c}^{*}$,
and $\lim_{t\to\infty}\xi(t)=\xi^{*}$. The fixed point values $x_{c}^{*}$,
$x^{*}$, $\xi_{c}^{*}$, and $\xi^{*}$ are all random variables.
It follows that the mean cavity abundance and the cavity correlation
become become time-independent random variables $\lim_{t\to\infty}\mu_{c}(t)=\mu_{c}^{*}=\langle x_{c}^{*}\rangle_{c}^{*}$,
and $\lim_{t,t'\to\infty}C_{c}(t,t')=q_{c}^{*}=\langle x_{c}^{*}x_{c}^{*}\rangle_{c}^{*}.$
Within this assumptions the cavity response function is time-translational-invariant,
i.e., $R_{c}(t,t')$ depends only on the time difference $\tau=t-t'$.
Due to causality $R_{c}(\tau)=0$ for $\tau\leq0$, and we define
$\chi_{c}^{*}=\int_{0}^{\infty}d\tau\,C_{c}^{*}(\tau)$. 

The fixed point satisfies the self-consistent equation 

\begin{equation}
x_{c}^{*}\left[1-x_{c}^{*}+\frac{K-1}{K}m\mu_{c}^{*}+\frac{K-1}{K}\left(\frac{m^{2}}{K}+\gamma\sigma^{2}\right)\chi_{c}^{*}x_{c}^{*}+\xi_{c}^{*}+h_{c}\right]=0,\label{eq:FP_eq_cavity}
\end{equation}

where $h_{c}$ is a constant external field $(h_{c}(t)=h_{c}$) needed
to obtain the integrated cavity response. The fixed point cavity noise is a Gaussian random variable with mean zero and variance $\Sigma_c^2q_c^*=(K-1)/K(m^2/K+\sigma^2)q_c^*$. We can therefore write is as $\xi_c^*=-\Sigma_c\sqrt{q_c^*}s$, where $s$ is a standard Gaussian random variable with zero mean and unit variance. Eq. \eqref{eq:FP_eq_cavity}
admits non-vanishing solution

\begin{equation}
x_{c}^{*}\left(s\right)=\frac{\Sigma_{c}\sqrt{q_{c}^{*}}(\Delta_c-s)+h_{c}}{1-\frac{K-1}{K}\left(\frac{m^{2}}{K}+\gamma\sigma^{2}\right)\chi_{c}^{*}}\Theta\left(\frac{\Sigma_{c}\sqrt{q_{c}^{*}}(\Delta_c-s)+h_{c}}{1-\frac{K-1}{K}\left(\frac{m^{2}}{K}+\gamma\sigma^{2}\right)\chi_{c}^{*}}\right),\label{eq:FP_sol_cavity}
\end{equation}

where $\Theta(x)$ is the Heaviside step function, which ensures positivity
of the fixed-point cavity abundance. For simplicity of notation we
have used the abbreviation $\Delta_{c}=(1+\frac{K-1}{K}m\mu_{c}^{*})/(\Sigma_{c}\sqrt{q_{c}^{*}})$. In a similar way, one finds the
fixed-point representative abundance

\begin{equation}
x^{*}\left(\xi^{*}\right)=\frac{1+m\mu_{c}^{*}+\xi^{*}}{1-\left(\frac{m^{2}}{K}+\gamma\sigma^{2}\right)\chi_{c}^{*}}\Theta\left(\frac{1+m\mu_{c}^{*}+\xi^{*}}{1-\left(\frac{m^{2}}{K}+\gamma\sigma^{2}\right)\chi_{c}^{*}}\right).\label{eq:FP_sol_marginal}
\end{equation}

We perform the average over the ensemble of fixed points following
the derivation of \cite{opperPhaseTransitionNoise1992,gallaRandomReplicatorsAsymmetric2006a,gallaDynamicallyEvolvedCommunity2018},
obtaining the integral equations (see Appendix \ref{sec:FP_RRG} for
a complete derivation)

\begin{subequations}
\begin{align}
\mu_{c}^{*} & =\frac{\Sigma_{c}\sqrt{q_{c}^{*}}}{1-\frac{K-1}{K}\left(\frac{m^{2}}{K}+\gamma\sigma^{2}\right)\chi_{c}^{*}}\int_{-\infty}^{\Delta_{c}}Ds\:\left(\Delta_{c}-s\right)\label{eq:FP_eq_1}\\
1 & =\frac{\Sigma_{c}^{2}}{\left[1-\frac{K-1}{K}\left(\frac{m^{2}}{K}+\gamma\sigma^{2}\right)\chi_{c}^{*}\right]^{2}}\int_{-\infty}^{\Delta_{c}}Ds\:\left(\Delta_{c}-s\right)^{2}\label{eq:FP_eq_2}\\
\chi_{c}^{*} & =\frac{1}{1-\frac{K-1}{K}\left(\frac{m^{2}}{K}+\gamma\sigma^{2}\right)\chi_{c}^{*}}\int_{-\infty}^{\Delta_{c}}Ds,\label{eq:FP_eq_3}
\end{align}
\end{subequations}

where $Ds=e^{-s^{2}/2}/\sqrt{2\pi}$. 

We recall that the integral equations are valid only in the stable
fixed point phase, and require $1-\frac{K-1}{K}\left(\frac{m^{2}}{K}+\gamma\sigma^{2}\right)\chi_{c}^{*}>0$.
Moreover, being $\mu_{c}^{*}$ the mean species abundance in the cavity
process, it must be non-negative. The integral equations can be solved
by Newton-Raphson methods, and deliver the cavity quantities $\mu_{c}^{*}$,
$q_{c}^{*}$ and $\chi_{c}^{*}$ as a function of the model parameters
$\sigma$, $m$ and $\gamma$. Once those quantities have been computed,
they can be inserted into Eq. \eqref{eq:FP_sol_marginal} to obtain
the fixed point distribution of the representative species abundance
$x^{*}$, which has a delta-peak at zero, representing the fraction
of species that go extinct, plus a clipped Gaussian, representing
the distribution of surviving species,

\begin{equation}
P(x^{*})=(1-\phi)\delta(x^{*})+p_{\text{surv}}(x^{*}).
\end{equation}

The distribution of surviving species is a clipped Gaussian

\begin{equation}
p_{\text{surv}}(x^{*})=\frac{1}{\sqrt{2\pi\Sigma_{\text{surv}}^{2}q_{c}^{*}}}e^{-\frac{(x^{*}-x_{S})^{2}}{2\Sigma_{\text{surv}}^{2}q_{c}^{*}}}\Theta(x^{*}),
\end{equation}

with 

\begin{equation}
x_{S}=\frac{1+m\mu_{c}^{*}}{1-\left(m^{2}/K+\gamma\sigma^{2}\right)\chi_{c}^{*}},\qquad\Sigma_{\text{surv}}=\frac{m^{2}/K+\sigma^{2}}{1-\left(m^{2}/K+\gamma\sigma^{2}\right)\chi_{c}^{*}}.
\end{equation}

The weight of the delta-peak at $x^{*}=0$ is the fraction of species
that go extinct $1-\phi$. Consequently the fraction of surviving
species is $\phi=\int_{0}^{\infty}dx^{*}\,p_{\text{surv}}(x^{*})=\left(1+\text{erf}(\Delta/\sqrt{2})\right)$
with $\Delta=\left(1+m\mu_{c}^{*}\right)/\left((m^{2}/K+\sigma^{2})\sqrt{q_{c}^{*}}\right)$.

\subsubsection{Stability analysis}

The fixed points we have studied in Section \ref{subsec:FP_analysis_RRG}
are defined only in the stable phase of the Lotka-Volterra system,
where the representative abundance converges to a unique stable fixed
point. We study the stability of those fixed points through a linear
stability analysis, in order to determine the boundaries of such a
stable phase in the phase diagram. A detailed derivation can be found
in Appendix \ref{sec:FP_stability_analysis_RRG}, where we adapted
the derivation of \cite{opperPhaseTransitionNoise1992,gallaDynamicallyEvolvedCommunity2018}.
Once the dynamics Eq. \eqref{eq:eff_proc_cav} reaches the non-vanishing
fixed point $x_{c}^{*}$, a small fluctuating perturbation $\varepsilon\zeta(t)$
is switched on, resulting in a fluctuation $y(t)$ about $x_{c}^{*}$
which evolves under the stochastic process

\begin{equation}
\frac{d}{dt}y(t)=x_{c}^{*}\left[-y(t)+\frac{K-1}{K}\left(\frac{m^{2}}{K}+\gamma\sigma^{2}\right)\int_{0}^{t}dt'\,R_{c}^{*}(t,t')y(t')+\nu(t)+\zeta(t)\right],\label{eq:stability_analysis_eq}
\end{equation}

where $\nu(t)$ is the deviation about the noise $\xi_{c}^{*}$. The
perturbing noise $\zeta(t)$ is Gaussian with zero mean and unit variance.
We can safely assume that $\nu(t)$ and $\zeta(t)$ are independent
from $x_{c}^{*}$ and $\xi_{c}^{*}$. Then one finds from self-consistency
$\langle\nu(t)\nu(t')\rangle=\frac{K-1}{K}(\frac{m^{2}}{K}+\sigma^{2})\langle y(t)y(t')\rangle$.
Passing to the Fourier space in Eq. \eqref{eq:stability_analysis_eq}

\begin{equation}
\tilde{y}(\omega)=\frac{\tilde{\nu}(\omega)+\tilde{\zeta}(\omega)}{\frac{{\rm i}\omega}{x_{c}^{*}}+\left(1-\frac{K-1}{K}\left(\frac{m^{2}}{K}+\gamma\sigma^{2}\right)\tilde{R_{c}^{*}}(\omega)\right)}.
\end{equation}

Focusing on the long-time behaviour of the correlation of fluctuations
we find

\begin{equation}
\langle\left|\tilde{y}(0)\right|^{2}\rangle=\frac{\phi_{c}}{\left[1-\frac{K-1}{K}\left(\frac{m^{2}}{K}+\gamma\sigma^{2}\right)\chi_{c}^{*}\right]^{2}-\frac{K-1}{K}\left(\frac{m^{2}}{K}+\sigma^{2}\right)\phi_{c}},
\end{equation}

where $\phi_{c}=\int_{-\infty}^{\Delta_{c}}Ds$ is the fraction of
surviving species in the cavity process. The correlation diverges
when $\frac{K-1}{K}\left(\frac{m^{2}}{K}+\sigma^{2}\right)\phi_{c}=\left[1-\frac{K-1}{K}\left(\frac{m^{2}}{K}+\gamma\sigma^{2}\right)\chi_{c}^{*}\right]^{2}$,
which leads to $\Delta_{c}=0$ and $\phi_{c}=1/2$ from Eq. \eqref{eq:FP_eq_2}.
From this we find that the system is in the stable phase for 

\begin{equation}
\frac{\left(2\frac{m^{2}}{K}+(1+\gamma)\sigma^{2}\right)^{2}}{\frac{m^{2}}{K}+\sigma^{2}}<2\frac{K}{K-1}.
\end{equation}

The result is consistent with the Fully Connected (FC) case, which
corresponds to taking the limit of infinite connectivity $K\to\infty$.
In that case the stable-fixed-point phase is independent of the value
of $m$, and is determined by $\sigma^{\text{FC}}<\sigma_{c}^{\text{FC}}$,
where $\left(\sigma_{c}^{\text{FC}}\right)^{2}=2/(1+\gamma)^{2}$.

We note that sparseness reduce the stability of the system when only
predator-prey interactions are present, i.e. when $\gamma=-1$. In
that limit the system is stable when $\left(\sigma^{\text{pp}}\right)^{2}>\left(\sigma_{c}^{\text{pp}}\right)^{2}=\frac{m^{2}}{K}\left(2\frac{K-1}{K^{2}}m^{2}-1\right)$,
which tends to zero in the FC limit, making the system stable in the
whole phase space.

\subsection{Proper derivation: population dynamics}
In the limit $T\to0$ we can assume that the system reaches a steady
state which does not depends on the initial condition. We further
assume that the dynamics reaches a unique fixed-point
$\lim_{t\to\infty}x_{i\setminus j}(t)=x_{i\setminus j}^{*}$, and therefore $\lim_{t\to\infty}\xi_{i\setminus j}(t)=\xi_{i\setminus j}^{*}$. The fixed point values $x_{i\setminus j}^{*}$,
$\xi_{i\setminus j}^{*}$, are all random variables. It follows that the mean cavity abundance and the cavity correlation
become become time-independent random variables $\lim_{t\to\infty}\mu_{i\setminus j}(t)=\mu_{i\setminus j}^{*}=\langle x_{i\setminus j}^{*}\rangle_{i\setminus j}^{*}$,
and $\lim_{t,t'\to\infty}C_{i\setminus j}(t,t')=q_{i\setminus j}^{*}=\langle ( x_{i\setminus j}^{*})^2\rangle_{i\setminus j}^{*}$.
Within this assumptions the cavity response function is time-translational-invariant,
i.e., $R_{i\setminus j}(t,t')$ depends only on the time difference $\tau=t-t'$.
Due to causality $R_{c}(i\setminus j)=0$ for $\tau\leq 0$, and we define
$\chi_{i\setminus j}^{*}=\int_{0}^{\infty}d\tau\,R_{i\setminus j}^{*}(\tau)$. 

The fixed point satisfies the self-consistent equation 

\begin{equation}
x_{i\setminus j}^{*}\left[1-x_{i\setminus j}^{*}+\sum_{k\in\partial i \setminus j}J_{ik}\mu_{k \setminus i}^{*}+\sum_{k\in\partial i \setminus j}J_{ik}J_{ki}\chi_{k \setminus i}^{*}x_{i \setminus j}^{*}+\xi_{i \setminus j}^{*}+h_{i \setminus j}\right]=0,\label{eq:FP_eq_cav_edge}
\end{equation}

where $h_{i \setminus j}$ is a constant external field needed
to obtain the integrated cavity response, and the noise is Gaussian with zero mean $\langle \xi_{i\setminus j}^*\rangle$ and variance $\langle (\xi_{i\setminus j})^*\rangle=\sum_{k\in\partial i\setminus j}J_{ik}^2q_{k\setminus i}^*=\Sigma_{i\setminus j}^2$. We can therefore write $\xi_{i\setminus j}^*=-\Sigma_{i\setminus j}s$, where $s$ is a standard Gaussian variable with zero mean and unit variance. Eq. \eqref{eq:FP_eq_cav_edge} admits non-vanishing solution

\begin{equation}
x_{i\setminus j}^{*}\left(s\right)=\frac{\Sigma_{i\setminus j}(\Delta_{i\setminus j}-s)+h_{i\setminus j}}{\Gamma_{i\setminus j}}\Theta\left(\frac{\Sigma_{i\setminus j}(\Delta_{i\setminus j}-s)+h_{i\setminus j}}{\Gamma_{i\setminus j}}\right),\label{eq:FP_sol_cavity}
\end{equation}

where $\Theta(x)$ is the Heaviside step function, which ensures positivity
of the fixed-point cavity abundance. We have defined $\Delta_{i\setminus j}=(1+\sum_{k\in\partial i \setminus j}J_{ik}\mu_{k \setminus i}^{*})/\Sigma_{i\setminus j}$, $\Gamma_{i\setminus j}=1-\sum_{k\in\partial i \setminus j}J_{ik}J_{ki}\chi_{k \setminus i}^{*}$ and we have assumed that $\sum_{k\in\partial i \setminus j}J_{ik}J_{ki}\chi_{k \setminus i}^{*}<1$.

By averaging over the Gaussian distribution of $s$ we obtain a set of self-consistent equations for the parameters

\begin{subequations}
\begin{align}
\mu_{i\setminus j}^*&=\frac{\Sigma_{i\setminus j}}{\Gamma_{i\setminus j}}\left[\Theta(\Gamma_{i\setminus j})\int_{-\infty}^{\Delta_{i\setminus j}}Ds\,(\Delta_{i\setminus j}-s)+\Theta(-\Gamma_{i\setminus j})\int_{\Delta_{i\setminus j}}^{\infty}Ds\,(\Delta_{i\setminus j}-s)\right]\\
q_{i\setminus j}^*&=\frac{\Sigma^2_{i\setminus j}}{\Gamma^2_{i\setminus j}}\left[\Theta(\Gamma_{i\setminus j})\int_{-\infty}^{\Delta_{i\setminus j}}Ds\,(\Delta_{i\setminus j}-s)^2+\Theta(-\Gamma_{i\setminus j})\int_{\Delta_{i\setminus j}}^{\infty}Ds\,(\Delta_{i\setminus j}-s)^2\right]\\
\chi_{i\setminus j}^*&=\frac{1}{\Gamma_{i\setminus j}}\left[\Theta(\Gamma_{i\setminus j})\int_{-\infty}^{\Delta_{i\setminus j}}Ds+\Theta(-\Gamma_{i\setminus j})\int_{\Delta_{i\setminus j}}^{\infty}Ds\right].
\end{align}
\end{subequations}

We observe that the integrals can be written in terms of the error function $\erf(x)=\int_0^{x} ds\,e^{-s^2}2/\sqrt{\pi}$. Defining $\varphi(x)=e^{-x^2/2}/\sqrt{2\pi}$ the normal distribution function and $\Phi(x)=[1+\erf(x/\sqrt{2})]/2$ its cumulative distribution function we obtain

\begin{subequations}
\begin{align}
\mu_{i\setminus j}^*&=\frac{\Sigma_{i\setminus j}}{\Gamma_{i\setminus j}}\left[\Theta(\Gamma_{i\setminus j})\left(\Delta_{i\setminus j}\Phi(\Delta_{i\setminus j})+\varphi(\Delta_{i\setminus j}) \right)+\Theta(-\Gamma_{i\setminus j}) \left(\Delta_{i\setminus j}\Phi(-\Delta_{i\setminus j})-\varphi(-\Delta_{i\setminus j}) \right)\right]\\
q_{i\setminus j}^*&=\frac{\Sigma^2_{i\setminus j}}{\Gamma^2_{i\setminus j}}\left[\Theta(\Gamma_{i\setminus j})\left(\left(1+\Delta_{i\setminus j}^2\right)\Phi(\Delta_{i\setminus j})+\Delta_{i\setminus j}\varphi(\Delta_{i\setminus j})\right)\right.\nonumber\\
&\quad \left.+\Theta(\Gamma_{i\setminus j})\left(\left(1+\Delta_{i\setminus j}^2\right)\Phi(-\Delta_{i\setminus j})-\Delta_{i\setminus j}\varphi(-\Delta_{i\setminus j})\right)  \right]\\
\chi_{i\setminus j}^*&=\frac{1}{\Gamma_{i\setminus j}}\left[\Theta(\Gamma_{i\setminus j})\Phi(\Delta_{i\setminus j})+\Theta(-\Gamma_{i\setminus j})\Phi(-\Delta_{i\setminus j})\right].
\end{align}
\end{subequations}

{\bf We can solve these equations by means of a population dynamics algorithm, here is their explicit version, together with the definitions of $\Delta_{i\setminus j}$ and $\Sigma_{i\setminus j}$
\begin{subequations}
\begin{align}
\Delta_{i\setminus j}&=\frac{1+\sum_{k\in \partial i\setminus j}J_{ik}\mu_{k\setminus i}^*}{\sqrt{\sum_{k\in \partial i\setminus j}J_{ik}^2 q_{k\setminus i}^*}}\\
\Sigma_{i\setminus j}^2&=\sum_{k\in \partial i\setminus j}J_{ik}^2 q_{k\setminus i}^*\\
\Gamma_{i\setminus j}&=1-\sum_{k\in\partial i \setminus j}J_{ik}J_{ki}\chi_{k \setminus i}^{*}.
\end{align}
\end{subequations}
}


\subsection{Proper derivation: just an idea...}

From the effective process of the cavity population density Eq.~\eqref{eq:eff_proc_cav}
one can obtain dynamical self-consistent equations for the averages
$\mu_{i\setminus j}$, $C_{i\setminus j}$, and $R_{i\setminus j}$.

\begin{subequations}
\begin{align}
\frac{d}{dt}\mu_{i\setminus j}(t) & =\langle\dot{x}_{i\setminus j}(t)\rangle_{i\setminus j}=\mu_{i\setminus j}(t)-C_{i\setminus j}(t,t)+\sum_{k\in\partial i\setminus j}J_{ik}\mu_{k\setminus i}(t)\mu_{i\setminus j}(t)\nonumber \\
 & \quad+\sum_{k\in\partial i\setminus j}\int_{0}^{t}dt'\,J_{ik}J_{ki}R_{k\setminus i}(t,t')C_{i\setminus j}(t,t')+\langle x_{i\setminus j}(t)\xi_{i\setminus j}(t)\rangle_{i\setminus j}\\
 & =\mu_{i\setminus j}(t)-C_{i\setminus j}(t,t)+\sum_{k\in\partial i\setminus j}J_{ik}\mu_{k\setminus i}(t)\mu_{i\setminus j}(t)\nonumber \\
& \quad+\sum_{k\in\partial i\setminus j}\int_{0}^{t}dt'\,\left[J_{ik}J_{ki}R_{k\setminus i}(t,t')C_{i\setminus j}(t,t')+J_{ik}^{2}C_{k\setminus i}(t,t')R_{i\setminus j}(t,t')\right],
\label{eq:righteqformu}
\end{align}
\end{subequations}

where we used

\begin{subequations}
\begin{align}
\langle x_{i\setminus j}(t)\xi_{i\setminus j}(t')\rangle_{i\setminus j} & =\frac{\delta}{\delta{\rm i}\hat{x}_{i\setminus j}(t')}\langle x_{i\setminus j}(t)\rangle_{i\setminus j} \\
 & =\langle x_{i\setminus j}(t)\int dt''\,\left[2T\delta(t'-t''){\rm i}\hat{x}_{i\setminus j}(t'')+\sum_{k\in\partial i\setminus j}J_{ik}^{2}C_{k\setminus i}(t',t''){\rm i}\hat{x}_{i\setminus j}(t'')\right]\rangle_{i\setminus j} \\
 & =2TR_{i\setminus j}(t,t')+\sum_{k\in\partial i\setminus j}\int_{0}^{t}dt''\,J_{ik}^{2}C_{k\setminus i}(t',t'')R_{i\setminus j}(t,t'').
\end{align}
\end{subequations}

Deriving the correlation function

\begin{subequations}
\begin{align}
\frac{\partial}{\partial t}C_{i\setminus j}(t,t') & =\langle\dot{x}_{i\setminus j}(t)x_{i\setminus j}(t')\rangle_{i\setminus j}=C_{i\setminus j}(t,t')-\langle x_{i\setminus j}^{2}(t)x_{i\setminus j}(t')\rangle_{i\setminus j}+\sum_{k\in\partial i\setminus j}J_{ik}\mu_{k\setminus i}(t)C_{i\setminus j}(t,t')\nonumber \\
 & \quad+\sum_{k\in\partial i\setminus j}\int_{0}^{t}dt''\:J_{ik}J_{ki}R_{k\setminus i}(t,t'')\langle x_{i\setminus j}(t)x_{i\setminus j}(t')x_{i\setminus j}(t'')\rangle_{i\setminus j}+\langle x_{i\setminus j}(t)x_{i\setminus j}(t')\xi_{i\setminus j}(t)\rangle_{i\setminus j}\\
 & =C_{i\setminus j}(t,t')-\langle x_{i\setminus j}^{2}(t)x_{i\setminus j}(t')\rangle_{i\setminus j}+\sum_{k\in\partial i\setminus j}J_{ik}\mu_{k\setminus i}(t)C_{i\setminus j}(t,t')+2T\langle x_{i\setminus j}(t)x_{i\setminus j}(t'){\rm i}\hat{x}_{i\setminus j}(t)\rangle_{i\setminus j}\nonumber \\
 & \quad+\sum_{k\in\partial i\setminus j}\int_{0}^{t}dt''\:[J_{ik}J_{ki}R_{k\setminus i}(t,t'')\langle x_{i\setminus j}(t)x_{i\setminus j}(t')x_{i\setminus j}(t'')\rangle_{i\setminus j}\nonumber \\
 & \quad+J_{ik}^{2}C_{k\setminus i}(t,t'')\langle x_{i\setminus j}(t)x_{i\setminus j}(t'){\rm i}\hat{x}_{i\setminus j}(t'')\rangle_{i\setminus j}],
\end{align}
\end{subequations}

and finally for the response
\begin{subequations}
\begin{align}
\frac{\partial}{\partial t}R_{i\setminus j}(t,t') & =\langle\dot{x}_{i\setminus j}(t){\rm i}\hat{x}_{i\setminus j}(t')\rangle_{i\setminus j}=R_{i\setminus j}(t,t')-\langle x_{i\setminus j}^{2}(t){\rm i}\hat{x}_{i\setminus j}(t')\rangle_{i\setminus j}+\sum_{k\in\partial i\setminus j}J_{ik}\mu_{k\setminus i}(t)R_{i\setminus j}(t,t')\nonumber \\
 & \quad+\sum_{k\in\partial i\setminus j}\int_{0}^{t}dt''\:J_{ik}J_{ki}R_{k\setminus i}(t,t'')\langle x_{i\setminus j}(t){\rm i}\hat{x}_{i\setminus j}(t')x_{i\setminus j}(t'')\rangle_{i\setminus j}+\langle x_{i\setminus j}(t){\rm i}\hat{x}_{i\setminus j}(t')\xi_{i\setminus j}(t)\rangle_{i\setminus j}\\
 & =R_{i\setminus j}(t,t')-\langle x_{i\setminus j}^{2}(t){\rm i}\hat{x}_{i\setminus j}(t')\rangle_{i\setminus j}+\sum_{k\in\partial i\setminus j}J_{ik}\mu_{k\setminus i}(t)R_{i\setminus j}(t,t')+2T\langle x_{i\setminus j}(t){\rm i}\hat{x}_{i\setminus j}(t'){\rm i}\hat{x}_{i\setminus j}(t)\rangle_{i\setminus j}\nonumber \\
 & \quad+\sum_{k\in\partial i\setminus j}\int_{0}^{t}dt''\:[J_{ik}J_{ki}R_{k\setminus i}(t,t'')\langle x_{i\setminus j}(t){\rm i}\hat{x}_{i\setminus j}(t')x_{i\setminus j}(t'')\rangle_{i\setminus j}\nonumber \\
 & \quad+J_{ik}^{2}C_{k\setminus i}(t,t'')\langle x_{i\setminus j}(t){\rm i}\hat{x}_{i\setminus j}(t'){\rm i}\hat{x}_{i\setminus j}(t'')\rangle_{i\setminus j}],
\end{align}
\end{subequations}


\subsubsection{Trying to approximate something}

Since we are interested in studying the statistics of the fixed-points of the system, we consider the long-time limit $t,t'\to\infty$. We can assume that the system retain no memory of the initial state distribution, therefore both the response and the correlation become Time Translational Invariant (TTI), i.e.  $R(t,t')=R(\tau=t-t')$ and $C(t,t')=C(\tau=t-t')$. Eq. \eqref{eq:righteqformu} becomes
\begin{align}
  \frac{d}{dt}\mu_{i\setminus j}(t)& =\mu_{i\setminus j}(t)-C_{i\setminus j}(0)+\sum_{k\in\partial i\setminus j}J_{ik}\mu_{k\setminus i}(t)\mu_{i\setminus j}(t)\nonumber \\
& \quad+\sum_{k\in\partial i\setminus j}\int d\tau\,\left[J_{ik}J_{ki}R_{k\setminus i}(\tau)C_{i\setminus j}(\tau)+J_{ik}^{2}C_{k\setminus i}(\tau)R_{i\setminus j}(\tau)\right].
\label{eq:righteqformuTTI}
\end{align}

Applying the Laplace transform to the expression above we obtain
\begin{align}
  z \hat{\mu}_{i\setminus j}(z) - \mu(0)& = \hat{\mu}_{i\setminus j}(z) - C_{i\setminus j}(0)\delta(z) + \sum_{k\in\partial i\setminus j}J_{ik} \mathcal{LT}\left[\mu_{k\setminus i}(t) {\mu}_{i\setminus j}(t)\right]
 \nonumber \\ 
& \quad+\frac{1}{z} \sum_{k\in\partial i\setminus j}\left[ J_{ik}J_{ki} \mathcal{LT}[ R_{k\setminus i}(\tau)C_{k\setminus i}(\tau) ] + J_{ik}^{2}\mathcal{LT}[ C_{k\setminus i}(\tau) R_{i\setminus j}(\tau)] \right].
\label{eq:equationmuLapExact}
\end{align}

Up to this point is essentially exact. Then, we decouple the Laplace transforms, assuming that $\mathcal{LT}[AB]\approx \mathcal{LT}[A]\mathcal{LT}[B]$. ({\bf Within this approximation we are closing our eyes})
\begin{align}
  z \hat{\mu}_{i\setminus j}(z) - \mu(0)& = \hat{\mu}_{i\setminus j}(z) - C_{i\setminus j}(0)\delta(z) + \sum_{k\in\partial i\setminus j}J_{ik}\hat{\mu}_{k\setminus i}(z) \hat{\mu}_{i\setminus j}(z)
\nonumber \\ 
& \quad+\frac{1}{z} \sum_{k\in\partial i\setminus j}\left[ J_{ik}J_{ki}  \hat{R}_{k\setminus i}(z)\hat{C}_{k\setminus i}(z) ] + J_{ik}^{2} \hat{C}_{k\setminus i}(z) \hat{R}_{i\setminus j}(z)]\right].
\label{eq:equationmuLapApprox}
\end{align}
 This leads to the self-consistent solution for $\mu$
\begin{equation}
  \hat{\mu}_{i\setminus j}(z) = \frac{\mu(0)-C_{i\setminus j}(0)\delta(z)+\frac{1}{z} \sum_{k\in\partial i\setminus j}\left[ J_{ik}J_{ki} \hat{R}_{k\setminus i}(z)\hat{C}_{k\setminus i}(z) + J_{ik}^{2} \hat{C}_{k\setminus i}(z) \hat{R}_{i\setminus j}(z) \right]}{ z - 1 -  \sum_{k\in\partial i\setminus j}J_{ik}\hat{\mu}_{k\setminus i}(z)}.
    \label{eq:FequationmuLapApprox}
\end{equation}
Notice that, for the Final Value Theorem (FVT), if $\mu_{i\setminus j}(t)$ is finite at infinite time, this finite value is defined by $\mu_{i\setminus j}^{eq} =\lim_{p \rightarrow 0} p \hat{\mu}_{i\setminus j}(p)$. Therefore
\begin{equation}
  \mu_{i\setminus j}^{eq} =   -\frac{\sum_{k\in\partial i\setminus j}\Big[ J_{ik}J_{ki} \hat{R}_{k\setminus i}(0)\hat{C}_{k\setminus i}(0) + J_{ik}^{2} \hat{C}_{k\setminus i}(0) \hat{R}_{i\setminus j}(0) \Big]}{\sum_{k\in\partial i\setminus j}J_{ik}\hat{\mu}_{k\setminus i}(0)+1}
    \label{eq:muatEquilibrium}
\end{equation}
{\bf Here we have a big problem, if we assume that $\lim_{z\to 0} z \hat{\mu}(z) = \mu^{eq}$ exists and is finite, then $\hat{\mu}(z) \approx \mu^{eq}/z$ for $z\to 0$, which means that $\hat{\mu}(0)$ diverges ... It follows that the denominator diverges ...}

We can do something similar for the remaining two quantities. Let us first focus on the response. First we must break down the correlation functions that involve more than two variables
\begin{subequations}
\begin{align}
\langle x_{i\setminus j}^{2}(t){\rm i}\hat{x}_{i\setminus j}(t')\rangle_{i\setminus j} & \approx \langle x_{i\setminus j}(t)\rangle \langle x_{i\setminus j}(t){\rm i}\hat{x}_{i\setminus j}(t')\rangle_{i\setminus j} =\mu_{i \setminus j}(t) R_{i \setminus j}(t,t') \\ 
\langle x_{i\setminus j}(t){\rm i}\hat{x}_{i\setminus j}(t'){\rm i}\hat{x}_{i\setminus j}(t)\rangle_{i\setminus j} &  \approx \langle x_{i\setminus j}(t)\rangle \langle {\rm i}\hat{x}_{i\setminus j}(t'){\rm i}\hat{x}_{i\setminus j}(t)\rangle_{i\setminus j} = 0 \\ 
    \langle x_{i\setminus j}(t){\rm i}\hat{x}_{i\setminus j}(t')x_{i\setminus j}(t'')\rangle_{i\setminus j}  & \approx \left[\langle x_{i\setminus j}(t)\rangle \langle {\rm i}\hat{x}_{i\setminus j}(t')x_{i\setminus j}(t'')\rangle_{i\setminus j} + \langle x_{i\setminus j}(t'')\rangle \langle {\rm i}\hat{x}_{i\setminus j}(t)x_{i\setminus j}(t')\rangle_{i\setminus j}\right]/2 \nonumber \\
    &=[\mu_{i\setminus j}(t) R_{i\setminus j}(t'',t')+\mu_{i\setminus j}(t'') R_{i\setminus j}(t,t')]/2
\end{align}
\end{subequations}

With this we can go further and rewrite the equation for $R$ in terms of one and two point correlation functions alone.


\begin{align}
  \frac{\partial}{\partial t}R_{i\setminus j}(t,t') & =R_{i\setminus j}(t,t')-
  \mu_{i \setminus j}(t) R_{i \setminus j}(t,t')
  +\sum_{k\in\partial i\setminus j}J_{ik}\mu_{k\setminus i}(t)R_{i\setminus j}(t,t')\nonumber\\
  &\quad + \mu_{i\setminus j}(t) \sum_{k\in\partial i\setminus j}\int_{0}^{t}dt''\:J_{ik}J_{ki}R_{k\setminus i}(t,t'')  R_{i\setminus j}(t',t'') \nonumber\\ 
  &= \Big[1 - \mu_{i \setminus j}(t) + \sum_{k\in\partial i\setminus j}J_{ik}\mu_{k\setminus i}(t) \Big] R_{i \setminus j}(t,t') + \mu_{i\setminus j}(t) \sum_{k\in\partial i\setminus j}\int_{0}^{t}dt''\:J_{ik}J_{ki}R_{k\setminus i}(t,t'')  R_{i\setminus j}(t',t'')
\end{align}

For times large enough one can substitute above $\mu_{i \setminus j}(t)=\mu_{eq}$, and use again the time translational invariance to show that:

\begin{equation}
  \hat{R}_{i\setminus j}(p) = \frac{R_0}{ p -  \Big[1 - \mu^{eq}_{i \setminus j} + \sum_{k\in\partial i\setminus j}J_{ik}\mu^{eq}_{k\setminus i} \Big] -  \mu^{eq}_{i \setminus j}  \sum_{k\in\partial i\setminus j}
    J_{ik}J_{ki} R_{k\setminus i}(p)}
  \label{eq:EquationforRLapApprox}
\end{equation}


I feel that eqs (\ref{eq:FequationmuLapApprox}), (\ref{eq:muatEquilibrium}) and
(\ref{eq:EquationforRLapApprox}) are already similar in spirit to what you got before for the linear problem. {\bf Please, check for a while the computations, and if they are not too crazy, we keep going}
    


\begin{acknowledgments}
We wish to acknowledge the support of ....
\end{acknowledgments}

\bibliography{GECaM_Ecosystems.bib}

\appendix

\section{\label{sec:path_int_RRG}Effective dynamics on a RRG}

The calculation is based on the principles of \cite{taraboloGaussianApproximationDynamic2024a}.
The dynamical partition function of the Lotka-Volterra system can
be written as

\begin{equation}
\text{Z=\ensuremath{\left\langle \int\mathcal{D}\vec{x}\prod_{i}p_{i}\left(x_{i}\left(0\right)\right)\delta\left(\text{equation of motion}\right)\right\rangle _{\vec{\eta}}},}
\end{equation}

where $\int\mathcal{D}\vec{x}=\prod_{i}\int\mathcal{D}x$ is a product
of functional integrals, $\langle\dots\rangle_{\vec{\eta}}=\left(\prod_{i}\int\mathcal{D}\eta_{i}\right)\dots$
is the average over the noise realizations, and $\delta(\dots)$ is
a functional Dirac delta function which restrict the integral to the
paths that satisfy the stochastic dynamics. Using the Fourier representation
of the delta function $\delta(x)\propto\int\mathcal{D}\hat{x}\,e^{-{\rm i}\int dt\,\hat{x}(t)x(t)}$
we obtain

\begin{equation}
Z=\left\langle \int\mathcal{D}\vec{x}\mathcal{D}\vec{\hat{x}}\prod_{i}p_{i}\left(x_{i}(0)\right)e^{\int dt\,({-{\rm i}}\hat{x}_{i}(t))\left(\frac{\dot{x_{i}}(t)}{x_{i}(t)}-1+x_{i}-\sum_{j\neq i}a_{ij}J_{ij}x_{j}-\eta_{i}-h_{i}(t)\right)}\right\rangle _{\vec{\eta}}.
\end{equation}

The integral over time run over a fixed time interval, but we omit
the limits of integration, here and elsewhere, for the sake of simplifying
the notation. Integrating over the Gaussian noise we obtain

\begin{align}
Z & =\int\mathcal{D}\vec{x}\mathcal{D}\vec{\hat{x}}\prod_{i}p_{i}\left(x_{i}(0)\right)e^{\int dt\,\left[{-{\rm i}}\hat{x}_{i}(t)\left(\frac{\dot{x_{i}}(t)}{x_{i}(t)}-1+x_{i}-h_{i}(t)\right)-T\left(\hat{x}_{i}(t)\right)^{2}\right]}\nonumber \\
 & \quad\times\prod_{i<j}e^{\int dt\,\left(a_{ij}J_{ij}{\rm i}\hat{x}_{i}(t)x_{j}(t)+a_{ji}J_{ji}{\rm i}\hat{x}_{j}(t)x_{i}(t)\right)}.
\end{align}

We put forward the following \textit{dynamic cavity equations} as
an ansatz for describing linearly coupled stochastic dynamics on sparse
networks

\begin{align}
c_{i\setminus j}\left[x_{i},\hat{x}_{i}\right] & =\frac{1}{Z_{ij}^{{\rm cav}}}e^{\int dt\,\left[{-{\rm i}}\hat{x}_{i}(t)\left(\frac{\dot{x_{i}}(t)}{x_{i}(t)}-1+x_{i}-h_{i}(t)\right)-T\left(\hat{x}_{i}(t)\right)^{2}\right]}\nonumber \\
 & \quad\times\prod_{k\in\partial i\setminus j}\int\mathcal{D}x_{k}\mathcal{D}\hat{x}_{k}c_{k\setminus i}\left[x_{k},\hat{x}_{k}\right]e^{\int dt\,\left(J_{ik}{\rm i}\hat{x}_{i}(t)x_{k}(t)+J_{ki}{\rm i}\hat{x}_{k}(t)x_{i}(t)\right)},\label{eq:cavity_marg_RRG}
\end{align}

with an arbitrary normalization factor 

\begin{align}
Z_{ij}^{{\rm cav}} & =\int\mathcal{D}x_{i}\mathcal{D}\hat{x}_{i}e^{\int dt\,\left[{-{\rm i}}\hat{x}_{i}(t)\left(\frac{\dot{x_{i}}(t)}{x_{i}(t)}-1+x_{i}-h_{i}(t)\right)-T\left(\hat{x}_{i}(t)\right)^{2}\right]}\nonumber \\
 & \quad\times\prod_{k\in\partial i\setminus j}\int\mathcal{D}x_{k}\mathcal{D}\hat{x}_{k}c_{k\setminus i}\left[x_{k},\hat{x}_{k}\right]e^{\int dt\,\left(J_{ik}{\rm i}\hat{x}_{i}(t)x_{k}(t)+J_{ki}{\rm i}\hat{x}_{k}(t)x_{i}(t)\right)}.
\end{align}

The full dynamic cavity marginals can be obtained from the cavity
ones as follows 

\begin{align}
c_{i}\left[x_{i},\hat{x}_{i}\right] & =\frac{1}{Z_{i}}e^{\int dt\,\left[{-{\rm i}}\hat{x}_{i}(t)\left(\frac{\dot{x_{i}}(t)}{x_{i}(t)}-1+x_{i}-h_{i}(t)\right)-T\left(\hat{x}_{i}(t)\right)^{2}\right]}\nonumber \\
 & \quad\times\prod_{k\in\partial i}\int\mathcal{D}x_{k}\mathcal{D}\hat{x}_{k}c_{k\setminus i}\left[x_{k},\hat{x}_{k}\right]e^{\int dt\,\left(J_{ik}{\rm i}\hat{x}_{i}(t)x_{k}(t)+J_{ki}{\rm i}\hat{x}_{k}(t)x_{i}(t)\right)},
\end{align}

and

\begin{align}
Z_{i} & =\int\mathcal{D}x_{i}\mathcal{D}\hat{x}_{i}e^{\int dt\,\left[{-{\rm i}}\hat{x}_{i}(t)\left(\frac{\dot{x_{i}}(t)}{x_{i}(t)}-1+x_{i}-h_{i}(t)\right)-T\left(\hat{x}_{i}(t)\right)^{2}\right]}\nonumber \\
 & \quad\times\prod_{k\in\partial i}\int\mathcal{D}x_{k}\mathcal{D}\hat{x}_{k}c_{k\setminus i}\left[x_{k},\hat{x}_{k}\right]e^{\int dt\,\left(J_{ik}{\rm i}\hat{x}_{i}(t)x_{k}(t)+J_{ki}{\rm i}\hat{x}_{k}(t)x_{i}(t)\right)}.
\end{align}

Expanding the cavity marginals Eq. \eqref{eq:cavity_marg_RRG} at
second order in the coupling parameters $J_{ij}$ we obtain a quadratic
form

\begin{align}
c_{i\setminus j}\left[x_{i},\hat{x}_{i}\right] & \propto\exp\left\{ \int dt\,\left[{-{\rm i}}\hat{x}_{i}(t)\left(\frac{\dot{x_{i}}(t)}{x_{i}(t)}-1+x_{i}-h_{i}(t)\right)-T\left(\hat{x}_{i}(t)\right)^{2}\right]+\sum_{k\in\partial i\setminus j}\int dt\,J_{ik}{\rm i}\hat{x}_{i}(t)\mu_{k\setminus i}(t)\right.\nonumber \\
 & \quad\left.+\frac{1}{2}\sum_{k\in\partial i\setminus j}\int dtdt'\,\left[J_{ik}^{2}{\rm i}\hat{x}_{i}(t)C_{k\setminus i}(t,t'){\rm i}\hat{x}_{i}(t')+2J_{ik}J_{ki}{\rm i}\hat{x}_{i}(t)R_{k\setminus i}(t,t')x_{i}(t')\right]\right\} ,
\end{align}

where $\mu_{k\setminus i}(t)$, $R_{k\setminus i}(t,t')$, and $C_{k\setminus i}(t,t')$
are respectively the cavity mean abundance, the cavity response function,
and the cavity (connected) two-point correlation function. They are
obtained self-consistently from the effective dynamics of the neighbors

\begin{equation}
\mu_{k\setminus i}(t)=\left\langle x_{k}(t)\right\rangle _{k\setminus i},\quad R_{k\setminus i}(t,t')=\left\langle x_{k}(t){\rm i}\hat{x}_{k}(t')\right\rangle _{k\setminus i},\quad C(t,t')=\left\langle x_{k}(t)x_{k}(t')\right\rangle _{k\setminus i},
\end{equation}

where $\left\langle \dots\right\rangle _{k\setminus i}=\int\mathcal{D}x_{k}\mathcal{D}\hat{x}_{k}\,\dots c_{k\setminus i}\left[x_{k},\hat{x}_{k}\right]$
describes the local average over the cavity marginal.

The argument of the exponential can be identified with a sort of local
effective single-edge action for the variables on site $i$ when the
interaction term with site $j$ is removed from the underlying graph.
Performing backward the steps of the path-integral representation,
we obtain an effective stochastic dynamics in terms of the stochastic
differential equation on the cavity graph

\begin{align}
\frac{d}{dt}x_{i\setminus j}(t) & =x_{i\setminus j}(t)\Biggl[1-x_{i\setminus j}(t)+\sum_{k\in\partial i\setminus j}J_{ik}\mu_{k\setminus i}(t)\nonumber \\
 & \quad+\sum_{k\in\partial i\setminus j}\int_{0}^{t}dt'J_{ik}R_{k\setminus i}(t,t')J_{ki}x_{i\setminus j}(t')+\xi_{i\setminus j}(t)\Biggr],
\end{align}

where $\xi_{i\setminus j}(t)$ is a Gaussian colored noise

\begin{equation}
\left\langle \xi_{i\setminus j}(t)\right\rangle _{i\setminus j}=0,\qquad\left\langle \xi_{i\setminus j}(t)\xi_{i\setminus j}(t')\right\rangle _{i\setminus j}=2T\delta(t-t')+\sum_{k\in\partial i\setminus j}J_{ik}^{2}C_{k\setminus i}(t,t').
\end{equation}

The terms $\mu_{k\setminus i}(t)$, $R_{k\setminus i}(t,t')$, and
$C_{k\setminus i}(t,t')$ have to be considered as external parameters,
obtained self-consistently as averages over the effective dynamics
of the neighbors. In a similar way, an effective process for the original
graph can be obtained as

\begin{align}
\frac{d}{dt}x_{i}(t) & =x_{i}(t)\Biggl[1-x_{i}(t)+\sum_{k\in\partial i}J_{ik}\mu_{k\setminus i}(t)\nonumber \\
 & \quad+\sum_{k\in\partial i}\int_{0}^{t}dt'J_{ik}R_{k\setminus i}(t,t')J_{ki}x_{i}(t')+\xi_{i}(t)+h_{i}(t)\Biggr],\label{eq:eff_proc_marg-1}
\end{align}

with Gaussian colored noise

\begin{equation}
\left\langle \xi_{i}(t)\right\rangle _{i}=0,\qquad\left\langle \xi_{i}(t)\xi_{i}(t')\right\rangle _{i}=2T\delta(t-t')+\sum_{k\in\partial i}J_{ik}^{2}C_{k\setminus i}(t,t').
\end{equation}


\section{Gaussian expansion of the disorder averaged dynamical partition function} \label{app:naive_method}

We focus here on the case of an heterogeneous random graph $G$, where
edges are randomly extracted from an adjacency matrix distribution
$p(A)$. A single instance $G$ of the random graph ensemble has a
set of fixed degrees for each node $k_{i}=\sum_{j}a_{ij}$, which
are sampled from a particular degree distribution $p(k)$. We denote
the average degree of the ensemble as $K=\sum_{k}kp(k)$. We can write
the distribution of the adjacency matrix as

\[
p(a_{ij},a_{ji})=\frac{K}{N}\delta_{a_{ij},1}\delta_{a_{ji},1}+\left(1-\frac{K}{N}\right)\delta_{a_{ij},0}\delta_{a_{ji},0}.
\]

We use the Gaussian Expansion Cavity method \cite{taraboloGaussianApproximationDynamic2024a}
to derive a Langevin equation for a representative species. This approach
begins by constructing the Martin-Siggia-Rose-Janssen-De Dominicis
(MSRJD) dynamic partition function through path integrals of the trajectories
\cite{martinStatisticalDynamicsClassical1973c,janssenLagrangeanClassicalField1976,dedominicisDynamicsSubstituteReplicas1978b}

\begin{subequations}
\begin{align}
Z & =\overline{\left\langle \int\mathcal{D}\vec{x}\prod_{i}\delta\left(\frac{\dot{x_{i}}(t)}{x_{i}(t)}-\left[1-x_{i}(t)+\sum_{j\neq i}a_{ij}J_{ij}x_{j}(t)+\eta_{i}(t)+h_{i}(t)\right]\right)\delta\left(\sum_{j}a_{ij}-k_{i}\right)\right\rangle _{\vec{\eta}}}\\
 & \propto\overline{\left\langle \int\mathcal{D}\vec{x}\mathcal{D}\vec{\hat{x}}d\vec{\psi}\prod_{i}e^{\int dt\,{\rm i}\hat{x}_{i}\left(\frac{\dot{x_{i}}}{x_{i}}-\left[1-x_{i}+\sum_{j\neq i}a_{ij}J_{ij}x_{j}+\eta_{i}+h_{i}\right]\right)}e^{{\rm i}\psi_{i}\left(\sum_{j}a_{ij}-k_{i}\right)}\right\rangle _{\vec{\eta}}}\\
 & \propto\int\mathcal{D}\vec{x}\mathcal{D}\vec{\hat{x}}d\vec{\psi}\prod_{i}e^{\int dt\,\left[{\rm i}\hat{x}_{i}\left(\frac{\dot{x_{i}}}{x_{i}}-1+x_{i}+h_{i}\right)-T\hat{x}_{i}^{2}\right]}e^{-{\rm i}\psi_{i}k_{i}}\overline{\prod_{i<j}e^{\left[a_{ij}\left({\rm i}\psi_{i}-\int dt\,J_{ij}{\rm i}\hat{x}_{i}x_{j}\right)+a_{ji}\left({\rm i}\psi_{j}-\int dt\,J_{ji}{\rm i}\hat{x}_{j}x_{i}\right)\right]}}\\
 & =\int\mathcal{D}\vec{x}\mathcal{D}\vec{\hat{x}}d\vec{\psi}\prod_{i}e^{\int dt\,\left[{\rm i}\hat{x}_{i}\left(\frac{\dot{x_{i}}}{x_{i}}-1+x_{i}+h_{i}\right)-T\hat{x}_{i}^{2}\right]}e^{-{\rm i}\psi_{i}k_{i}}\prod_{i<j}\left[1-\frac{K}{N}+\frac{K}{N}\overline{e^{\left[\left({\rm i}\psi_{i}-\int dt\,J{\rm i}\hat{x}_{i}x_{j}\right)+\left({\rm i}\psi_{j}-\int dt\,J'{\rm i}\hat{x}_{j}x_{i}\right)\right]}}\right]\\
 & \approx\int\mathcal{D}\vec{x}\mathcal{D}\vec{\hat{x}}d\vec{\psi}\prod_{i}e^{\int dt\,\left[{\rm i}\hat{x}_{i}\left(\frac{\dot{x_{i}}}{x_{i}}-1+x_{i}+h_{i}\right)-T\hat{x}_{i}^{2}\right]}e^{-{\rm i}\psi_{i}k_{i}}e^{-\frac{KN}{2}+\frac{K}{2N}\sum_{ij}\overline{e^{-\int dt\,\left(J{\rm i}\hat{x}_{i}x_{j}+J'{\rm i}\hat{x}_{j}x_{i}\right)}}e^{{\rm i}(\psi_{i}+\psi_{j})}}.
\end{align}
\end{subequations}

We introduce the usual functional order parameter following the steps
of \cite{aguirre-lopezHeterogeneousMeanfieldAnalysis2024a}

\begin{equation}
c[x,\hat{x}]=\frac{1}{N}\sum_{i}e^{-{\rm i}\psi_{i}}\delta(x-x_{i})\delta(\hat{x}-\hat{x}_{i}).
\end{equation}

This definition can be inserted into $Z$ using as usual functional
Dirac deltas
\begin{subequations}
\begin{align}
Z & \propto\int\mathcal{D}\vec{x}\mathcal{D}\vec{\hat{x}}d\vec{\psi}\prod_{i}e^{\int dt\,\left[{\rm i}\hat{x}_{i}\left(\frac{\dot{x_{i}}}{x_{i}}-1+x_{i}+h_{i}\right)-T\hat{x}_{i}^{2}\right]}e^{-{\rm i}\psi_{i}k_{i}}e^{\frac{K}{2N}\sum_{ij}\overline{e^{-\int dt\,\left(J{\rm i}\hat{x}_{i}x_{j}+J'{\rm i}\hat{x}_{j}x_{i}\right)}}e^{{\rm i}(\psi_{i}+\psi_{j})}}\nonumber \\
 & \quad\times\int\mathcal{D}c\,\delta\left(Nc[x,\hat{x}]-\sum_{i}e^{-{\rm i}\psi_{i}}\delta(x-x_{i})\delta(\hat{x}-\hat{x}_{i})\right)\\
 & =\int\mathcal{D}\vec{x}\mathcal{D}\vec{\hat{x}}d\vec{\psi}\prod_{i}e^{\int dt\,\left[{\rm i}\hat{x}_{i}\left(\frac{\dot{x_{i}}}{x_{i}}-1+x_{i}+h_{i}\right)-T\hat{x}_{i}^{2}\right]}e^{-{\rm i}\psi_{i}k_{i}}e^{\frac{K}{2N}\sum_{ij}\overline{e^{-\int dt\,\left(J{\rm i}\hat{x}_{i}x_{j}+J'{\rm i}\hat{x}_{j}x_{i}\right)}}e^{{\rm i}(\psi_{i}+\psi_{j})}}\nonumber \\
 & \quad\times\int\mathcal{D}c\mathcal{D}\hat{c}\,e^{{\rm i}\int\mathcal{D}x\mathcal{D}\hat{x}\,\hat{c}[x,\hat{x}]\left(Nc[x,\hat{x}]-\sum_{i}e^{-{\rm i}\psi_{i}}\delta(x-x_{i})\delta(\hat{x}-\hat{x}_{i})\right)}\\
 & =\int\mathcal{D}c\mathcal{D}\hat{c}\,e^{{\rm i}N\int\mathcal{D}x\mathcal{D}\hat{x}\,\hat{c}[x,\hat{x}]c[x,\hat{x}]+\frac{KN}{2}\int\mathcal{D}x\mathcal{D}\hat{x}\,c[x,\hat{x}]\int\mathcal{D}y\mathcal{D}\hat{y}\,c[y,\hat{y}]\overline{e^{-\int dt\,\left({\rm i}\hat{x}Jy+{\rm i}\hat{y}J'x\right)}}}\nonumber \\
 & \quad\times\prod_{k}\left[\int\mathcal{D}x\mathcal{D}\hat{x}e^{\int dt\,\left[{\rm i}\hat{x}\left(\frac{\dot{x}}{x}-1+x+h\right)-T\hat{x}^{2}\right]}\frac{\left(-{\rm i}\hat{c}\left[x,\hat{x}\right]\right)^{k}}{k!}\right]^{N_{k}} & .
\end{align}
\end{subequations}

We end up with the usual heterogeneous dynamical mean field theory,
already done by \cite{aguirre-lopezHeterogeneousMeanfieldAnalysis2024a}
.

\begin{equation}
Z\propto\int\mathcal{D}c\mathcal{D}\hat{c}\,e^{N\Psi[c,\hat{c}]}\sim e^{N\Psi^{*}},
\end{equation}

where

\begin{equation}
\Psi^{*}=\max_{c,\hat{c}}\left\{ \Xi_{1}[c,\hat{c}]+\Xi_{2}[c]+\Xi_{3}[\hat{c}]\right\} ,
\end{equation}

and
\begin{subequations}
\begin{align}
\Xi_{1}[c,\hat{c}] & ={\rm i}\int\mathcal{D}x\mathcal{D}\hat{x}\,\hat{c}[x,\hat{x}]c[x,\hat{x}]\\
\Xi_{2}[c] & =\frac{K}{2}\int\mathcal{D}x\mathcal{D}\hat{x}\,c[x,\hat{x}]\int\mathcal{D}y\mathcal{D}\hat{y}\,c[y,\hat{y}]\overline{e^{-\int dt\,\left({\rm i}\hat{x}Jy+{\rm i}\hat{y}J'x\right)}}\\
\Xi_{3}[\hat{c}] & =\frac{1}{N}\log\prod_{k}\left[\int\mathcal{D}x\mathcal{D}\hat{x}e^{\int dt\,\left[{\rm i}\hat{x}\left(\frac{\dot{x}}{x}-1+x+h\right)-T\hat{x}^{2}\right]}\frac{\left(-{\rm i}\hat{c}\left[x,\hat{x}\right]\right)^{k}}{k!}\right]^{N_{k}}\nonumber \\
 & \sim\sum_{k}p(k)\log\int\mathcal{D}x\mathcal{D}\hat{x}e^{\int dt\,\left[{\rm i}\hat{x}\left(\frac{\dot{x}}{x}-1+x+h\right)-T\hat{x}^{2}\right]}\frac{\left(-{\rm i}\hat{c}\left[x,\hat{x}\right]\right)^{k}}{k!}.
\end{align}
\end{subequations}

The saddle-point solution gives

\begin{subequations}
\begin{align}
\hat{c}[x,\hat{x}] & ={\rm i}K\int\mathcal{D}y\mathcal{D}\hat{y}\,c[y,\hat{y}]\overline{e^{-\int dt\,\left({\rm i}\hat{x}Jy+{\rm i}\hat{y}J'x\right)}}\\
c[x,\hat{x}] & =K\sum_{k}\frac{kp(k)}{K}\frac{e^{\int dt\,\left[{\rm i}\hat{x}\left(\frac{\dot{x}}{x}-1+x+h\right)-T\hat{x}^{2}\right]}\left(-{\rm i}\hat{c}\left[x,\hat{x}\right]\right)^{k-1}}{\int\mathcal{D}x\mathcal{D}\hat{x}\,e^{\int dt\,\left[{\rm i}\hat{x}\left(\frac{\dot{x}}{x}-1+x+h\right)-T\hat{x}^{2}\right]}\left(-{\rm i}\hat{c}\left[x,\hat{x}\right]\right)^{k}}.
\end{align}
\end{subequations}

Eliminating $\hat{c}$ leads to the following self-consistent equation
for the functional order parameter $c$

\begin{equation}
c[x,\hat{x}]=\sum_{k}\frac{kp(k)}{K}\frac{1}{Z_{\text{eff},k}}e^{-S_{\text{eff},k-1}[x,\hat{x}]},
\end{equation}

with the effective local action

\begin{subequations}
\begin{align}
S_{\text{eff},k}[x,\hat{x}] & =-\int dt\,\left[{\rm i}\hat{x}\left(\frac{\dot{x}}{x}-1+x+h\right)-T\hat{x}^{2}\right]\nonumber \\
 & \quad-k\log\int\mathcal{D}y\mathcal{D}\hat{y}\,c[y,\hat{y}]\overline{e^{-\int dt\,\left({\rm i}\hat{x}Jy+{\rm i}\hat{y}J'x\right)}}\\
Z_{\text{eff},k} & =\int\mathcal{D}x\mathcal{D}\hat{x}\,e^{-S_{\text{eff},k}[x,\hat{x}]}=1.
\end{align}
\end{subequations}

The partition function ensures the desired normalization of the functional
order parameter

\begin{equation}
\int\mathcal{D}x\mathcal{D}\hat{x}\,c[x,\hat{x}]=1.
\end{equation}

Expanding up to second order in $J$ the effective local action

\begin{subequations}
\begin{align}
S_{\text{eff},k}[x,\hat{x}] & \approx-\int dt\,\left[{\rm i}\hat{x}\left(\frac{\dot{x}}{x}-1+x+h\right)-T\hat{x}^{2}\right]+{\rm i}k\overline{J}\int dt\:\left(\hat{x}\langle y\rangle+\langle\hat{y}\rangle x\right)\nonumber \\
 & \quad-\frac{k}{2}\int dtdt'\,\left(\overline{J^{2}}\hat{x}(t)\langle y(t)y(t')\rangle\hat{x}(t')+2\overline{JJ'}\hat{x}(t)\langle y(t)\hat{y}(t')\rangle x(t')+\overline{J'{}^{2}}x(t)\langle\hat{y}(t)\hat{y}(t')\rangle x(t')\right)\\
 & =-\int dt\,\left[{\rm i}\hat{x}\left(\frac{\dot{x}}{x}-1+x+h\right)-T\hat{x}^{2}\right]+{\rm i}k\overline{J}\int dt\:\hat{x}(t)M_{\text{cav}}(t)\nonumber \\
 & \quad-\frac{k}{2}\int dtdt'\,\left(\overline{J^{2}}\hat{x}(t)C_{\text{cav}}(t,t')\hat{x}(t')+2\overline{JJ'}\hat{x}(t)R_{\text{cav}}(t,t')x(t')\right)
\end{align}
\end{subequations}

\section{\label{sec:FP_RRG}Fixed point derivation on a RRG}

The cavity fixed point solution is given by \eqref{eq:FP_sol_cavity}.
Since $\xi_{c}^{*}$ is Gaussian with zero average and variance $\langle\left(\xi_{c}^{*}\right)^{2}\rangle=\frac{K-1}{K}\left(\frac{m^{2}}{K}+\sigma^{2}\right)q_{c}^{*}=\Sigma_{c}^{2}q_{c}^{*}$
we can write $\xi_{c}^{*}=-\Sigma_{c}\sqrt{q_{c}^{*}}s$ where $s$
is a standard Gaussian random variable with zero mean and unit variance.
Defining $\Delta_{c}=\left(1-\frac{K-1}{K}m\mu_{c}^{*}\right)/\left(\Sigma_{c}\sqrt{q_{c}^{*}}\right)$
and assuming $1-\frac{K-1}{K}(\frac{m^{2}}{K}+\gamma\sigma^{2})\chi_{c}^{*}>0$
we can write

\begin{equation}
x_{c}^{*}(s)=\frac{\Sigma_{c}\sqrt{q_{c}^{*}}\left(\Delta_{c}-s\right)+h_{c}}{1-\frac{K-1}{K}(\frac{m^{2}}{K}+\gamma\sigma^{2})\chi_{c}^{*}}\Theta\left(\Sigma_{c}\sqrt{q_{c}^{*}}\left(\Delta_{c}-s\right)+h_{c}\right).
\end{equation}

Averaging over the distribution of $s$at $h_{c}=0$

\begin{subequations}
\begin{align}
\mu_{c}^{*} & =\langle x_{c}^{*}(s)\rangle_{s}=\frac{\Sigma_{c}\sqrt{q_{c}^{*}}}{1-\frac{K-1}{K}\left(\frac{m^{2}}{K}+\gamma\sigma^{2}\right)\chi_{c}^{*}}\int_{-\infty}^{\Delta_{c}}Ds\:\left(\Delta_{c}-s\right)\\
q_{c}^{*} & =\langle\left(x_{c}^{*}(s)\right)^{2}\rangle_{s}=\frac{\Sigma_{c}^{2}q_{c}^{*}}{\left[1-\frac{K-1}{K}\left(\frac{m^{2}}{K}+\gamma\sigma^{2}\right)\chi_{c}^{*}\right]^{2}}\int_{-\infty}^{\Delta_{c}}Ds\:\left(\Delta_{c}-s\right)^{2},
\end{align}
\end{subequations}

where $Ds=e^{-s^{2}/2}/\sqrt{2\pi}$. The integrated response $\chi_{c}^{*}$
is obtained as derivative of the average cavity population size with
respect to a constant vanishing field, acting at all times, i.e. $h_{c}(t)=h_{c}$

\begin{align}
\chi_{c}^{*} & =\left.\frac{d}{dh_{c}}\langle x_{c}^{*}(s)\rangle_{s}\right|_{h_{c}=0}=\frac{d}{dh_{c}}\frac{1}{1-\frac{K-1}{K}\left(\frac{m^{2}}{K}+\gamma\sigma^{2}\right)\chi_{c}^{*}}\int_{-\infty}^{\Delta_{c}+h_{c}/(\Sigma_{c}\sqrt{q_{c}^{*}})}Ds\:\left[\Sigma_{c}\sqrt{q_{c}^{*}}\left(\Delta_{c}-s\right)+h_{c}\right]\\
 & =\frac{1}{1-\frac{K-1}{K}\left(\frac{m^{2}}{K}+\gamma\sigma^{2}\right)\chi_{c}^{*}}\int_{-\infty}^{\Delta_{c}}Ds.
\end{align}


\section{\label{sec:FP_stability_analysis_RRG}Stability analysis
on a RRG}

In order to study the stability of the fixed points, we perform a
linear stability analysis. We add white noise $\zeta(t)$ of unit
variance to the cavity effective process

\begin{align}
\frac{d}{dt}x_{c}(t) & =x_{c}(t)\Biggl[1-x_{c}(t)+\frac{K-1}{K}m\mu_{c}(t)\nonumber \\
 & \quad+\frac{K-1}{K}\left(\frac{m^{2}}{K}+\gamma\sigma^{2}\right)\int_{0}^{t}dt'R_{c}(t,t')x_{c}(t')+\xi_{c}(t)+h_{c}(t)+\varepsilon\zeta(t)\Biggr].\label{eq:perturbed_cavity_process_RRG}
\end{align}

We study fluctuations $y(t)$ about a fixed point of Eq. \eqref{eq:eff_proc_cav_RRG},
i.e. we write $x_{c}(t)=x_{c}^{*}+\varepsilon y(t)$, and denote the
resulting additional term in the self-consistent noise by $\varepsilon\nu(t)$
(i.e., $\xi_{c}(t)=\xi_{c}^{*}+\varepsilon\nu(t)$). We use $\langle\zeta(t)\rangle=\langle\nu(t)\rangle=\langle y(t)\rangle=0$. 

\subsubsection{Extinct species}

We consider firstly the stability of the fixed point $x_{c}^{*}=0$.
Linearizing Eq. \eqref{eq:perturbed_cavity_process_RRG} in $\varepsilon$

\begin{equation}
\frac{d}{dt}y(t)=y(t)\left(1+\frac{K-1}{K}m\mu_{c}^{*}+\xi_{c}^{*}\right).
\end{equation}

Given our assumptions, the expression within the square brackets is
negative for fixed points at zero {[}refer to Eq. \eqref{eq:FP_sol_cavity}{]}.
This implies that any disturbances around these zero fixed points
diminish exponentially over time. Conversely, if the expression in
the square brackets is positive, the zero fixed point becomes unstable,
validating our choice to instead use the non-zero solution $x_{c}^{*}$
in such cases.

\subsubsection{Extant species}

We now focus on extant species, i.e. non-vanishing fixed points $x_{c}^{*}>0$.
Linearizing Eq. \eqref{eq:perturbed_cavity_process_RRG} in $\varepsilon$

\begin{equation}
\frac{d}{dt}y(t)=x_{c}^{*}\left[-y(t)+\frac{K-1}{K}\left(\frac{m^{2}}{K}+\gamma\sigma^{2}\right)\int_{0}^{t}dt'\,R_{c}^{*}(t,t')y(t')+\nu(t)+\zeta(t)\right],
\end{equation}

where $\langle\nu(t)\nu(t')\rangle=\frac{K-1}{K}(\frac{m^{2}}{K}+\sigma^{2})\langle y(t)y(t')\rangle$
by self-consistency. Passing to Fourier space we obtain

\begin{equation}
\frac{{\rm i}\omega}{x_{c}^{*}}\tilde{y}(\omega)=\left[\frac{K-1}{K}(\frac{m^{2}}{K}+\gamma\sigma^{2})\tilde{R}_{c}^{*}(\omega)-1\right]\tilde{y}(\omega)+\tilde{\nu}(\omega)+\tilde{\zeta}(\omega),
\end{equation}

from which we can write

\[
\tilde{y}(\omega)=\frac{\tilde{\nu}(\omega)+\tilde{\zeta}(\omega)}{\frac{{\rm i}\omega}{x_{c}^{*}}+\left(1-\frac{K-1}{K}\left(\frac{m^{2}}{K}+\gamma\sigma^{2}\right)\tilde{R_{c}^{*}}(\omega)\right)}.
\]

Noting that $\langle\left|\tilde{y}(\omega)\right|^{2}\rangle$ is
the Fourier transform of the correlation of the fluctuations $\langle y(t)y(t+\tau)\rangle$
in the stationary state where it only depends on $\tau$, we can study
its long-time behaviour at $\omega=0$. If this quantity diverges,
perturbations do not decay to zero, signalling an instability of the
fixed point. Using $\tilde{R_{c}^{*}}(\omega)=\int d\tau\:R_{c}^{*}(\tau)=\chi_{c}^{*}$

\begin{equation}
\langle\left|\tilde{y}(0)\right|^{2}\rangle=\frac{\phi_{c}}{\left[1-\frac{K-1}{K}\left(\frac{m^{2}}{K}+\gamma\sigma^{2}\right)\chi_{c}^{*}\right]^{2}-\frac{K-1}{K}\left(\frac{m^{2}}{K}+\sigma^{2}\right)\phi_{c}},
\end{equation}

where $\phi_{c}=\int_{-\infty}^{\Delta_{c}}Ds$, fraction of surviving
species in the cavity process, appears because we are only considering
non-zero fixed points (fluctuations about zero fixed points decay,
as demonstrated above, and so they do not contribute to $\langle\left|\tilde{y}(0)\right|^{2}\rangle$. 

The correlation diverges when $\frac{K-1}{K}\left(\frac{m^{2}}{K}+\sigma^{2}\right)\phi_{c}=\left[1-\frac{K-1}{K}\left(\frac{m^{2}}{K}+\gamma\sigma^{2}\right)\chi_{c}^{*}\right]^{2}$.
Inserting the critical condition into Eq. \eqref{eq:FP_eq_2} we find

\begin{equation}
\phi_{c}=\int_{-\infty}^{\Delta_{c}}Ds\left(\Delta_{c}-s\right)^{2}=\int_{-\infty}^{\Delta_{c}}Ds,
\end{equation}

from which $\Delta_{c}=0$ and $\phi_{c}=1/2$. Inserting these values
into Eqs. \eqref{eq:FP_eq_2} and \eqref{eq:FP_eq_3} and eliminating
$\chi_{c}^{*}$ we obtain the critical condition

\begin{equation}
\frac{\left(2\frac{m^{2}}{K}+(1+\gamma)\sigma^{2}\right)^{2}}{\frac{m^{2}}{K}+\sigma^{2}}=2\frac{K}{K-1}.
\end{equation}


\end{document}
